% §1 Introduction (lead). Sources: docs/01 §1 and §6, docs/02 §2, docs/11 §2 and §11.
% Every literature number below is context from the cited work, not a comparison.
\section{Introduction}
\label{sec:intro}

When a microservice system degrades, an on-call engineer has to work out which service, resource or change caused it, from metrics, logs, traces and change events that usually live in different tools.
One vendor survey puts the median cost of a high-impact outage at USD~2M per hour~\cite{newrelic2025forecast}; every cost figure we found is vendor-run or vendor-sponsored.
Most outages follow a change: in Google's 2010--2017 postmortems, binary and configuration pushes triggered 68\% of outages~\cite{googlesre2017postmortem}.
Yet the public benchmarks on which research methods are scored, such as RCAEval, PetShop and OpenRCA, carry no deploy, configuration or feature-flag events~\cite{rcaeval2024,petshop2024,openrca2025}.

Root-cause analysis (RCA) research has moved quickly towards large language model (LLM) agents.
On one 90-case benchmark, simple harnesses around frontier cloud models reach a mean PR@1 of 72.8\% (an SQL agent) and 77.2\% (a Codex harness) over three models, against 38.9\% for the best non-LLM method~\cite{rcatlas2026}.
Stricter tasks tell a different story.
Across eleven LLMs, the exact root-cause set is recovered in 20.7\% of OpenRCA~2.0 cases on average~\cite{openrca2026v2}; the best model solves 25.3\% of medium and 10.0\% of hard ORCA-bench tasks~\cite{orcabench2026}; and on SREGym one agent's end-to-end success falls from 63.7\% on problems ported from earlier benchmarks to 17.9\% on failures new to SREGym~\cite{sregym2026}.
Agents can also be steered: one planted false sentence made models name a decoy service first in 21--42\% of flagged RCAEval RE2 cases~\cite{hussain2026prereg}.
Explicit verification helps: an LLM-based verifier and grounder raised Acc@1 by 3.5--35.1 points in all 14 configurations of one study, but did not reduce errors in which the model misread the meaning of the telemetry~\cite{diagguard2026}.
We found no paper that isolates a purely deterministic verifier over LLM proposals.

The field's own measurements are fragile as well.
The same method on the same cases scores differently in different harnesses: BARO's AC@1 on RE2-TT is 0.67 in RCAEval's harness and 0.32 in PRISM's~\cite{rcaeval2024,prism2026}.
Shifting the assumed fault onset two minutes late dropped N-Sigma from 0.80 to 0.07 AC@1 on RE2~\cite{hussain2026prereg}.
Simple alert counting matches or beats state-of-the-art methods on Nezha, Eadro and RCAEval~\cite{fang2026rethinking}.
Pooled leaderboards pick the worse method on up to 5 of 11 subsystems, and none of 23 audited RCA papers released per-case outputs~\cite{hu2026audit}.

Three properties that an engineer needs before acting on an automated diagnosis are therefore missing together.
First, an LLM should be able to contribute without being trusted: its claims should count only after code has checked them, and it should not be able to argue for its favourite candidate.
Second, the system should say when it does not know: every method evaluated on PetShop named false root causes during normal operation~\cite{petshop2024}, and a recent LLM benchmark scores abstention as a wrong answer~\cite{rcatlas2026}.
Third, a verdict should be reproducible by someone else, which LLM inference does not give for free: 1,000 temperature-zero runs of one prompt produced 80 distinct outputs in one report~\cite{thinkingmachinesnondeterminism}.
Put together, we found no open, laptop-runnable RCA engine that ranks candidates with cheap deterministic methods, lets an LLM contribute without being trusted, says when it does not know, lets anyone replay a verdict, and publishes per-case accuracy on public benchmarks against honest baselines.

We present \sys, an RCA engine built on one rule: \emph{AI proposes, code decides}.
Six deterministic rankers and a local open-weight LLM of at most 10B parameters propose hypotheses; a deterministic verifier decides.
The LLM never asserts evidence.
It may only request named predicates from a fixed library, and code evaluates each requested predicate for \emph{every} candidate in the pool, so a check the LLM asks for can help whichever candidate it actually supports.
The verifier runs ten fixed trials on each candidate, combines them with the code-checked claims in a fixed scoring rule, and returns a split-conformal candidate set together with a status that may be \st{INCONCLUSIVE} or \NOFAULT.
Each verdict is sealed in a receipt from which anyone can replay it to the same hash.
The whole pipeline runs on a 16\,GB laptop without a GPU or a cloud service.

This paper makes four contributions.
\begin{itemize}
\item \textbf{C1, a propose--verify architecture with symmetric claim evaluation and replayable receipts} (\S\ref{sec:design}). Rankers and an LLM propose; a fixed decision function over code-recomputed predicates decides; every predicate the LLM requests is evaluated for every candidate; and verdicts replay from a receipt and a recording of the LLM calls.
\item \textbf{C2, conformal, set-valued RCA with a no-fault answer and abstention} (\S\ref{sec:conformal}). To our knowledge, based on searches up to October 2026, no microservice RCA work outputs conformal candidate sets with a no-fault option. Close work calibrates the confidence of LLM root-cause answers~\cite{lmpace2024}, gives distribution-free confidence sets for which stream changed first~\cite{conch2026}, or abstains in RCA for network fabrics~\cite{abstention2026osdi}. We report coverage, set size and risk--coverage curves.
\leadq{docs/11 §11: CONCH/CROC and the OSDI'26 abstention paper were seen only as snippets; read both in full before keeping this sentence.}
\item \textbf{C3, a decomposed, onset-robust, per-system evaluation with released per-case outputs} (\S\ref{sec:eval}). We separate retrieval from re-ranking~\cite{decomprca2026}, report every result per system next to chance, a prior-only ranker, N-Sigma, BARO and an alert-count baseline, test with clustered paired statistics, compare oracle and detected alert times, and froze all of it as the git tag \code{prereg-v1} before the first test-split run of any \sys arm.
\item \textbf{C4, a change-aware evidence channel stress-tested with decoy changes} (\S\ref{sec:rq4}). Change proximity and changed-versus-peer contrast enter as evidence, and RCAEval-C adds a true change with probability~$\pi$ plus decoy changes to RCAEval RE2. Its change logs are synthetic, so we report it as a test of how gullible the channel is, not as evidence of real-world gains.
\end{itemize}

% Headline results: neutral sentences until bench/results/ exists; the wording
% ("not worse", "better") is chosen only after the numbers are known.
\leadq{headline results, one sentence per RQ, written at M4 from the generated numbers.}
Pooled over RE2, the paired AC@1 difference between the AI-off pipeline and the best reproduced simple baseline is \resnum{re2-oracle-rq1-v1}{pooled AC@1 difference, veyltrace:off minus best simple baseline} (95\% CI \resnum{re2-oracle-rq1-v1}{pooled AC@1 difference CI}).
Code-verified proposals change AC@1 by \resnum{re2-oracle-rq2-v1}{pooled proposer lift, veyltrace:brief minus veyltrace:off}, and planted decoy proposals change the top-1 candidate in \resnum{re2-oracle-planted-v1}{decoy_top1_change} of cases.
Held-out coverage of the candidate sets at a 90\% target is \resnum{re2-oracle-rq3-v1}{coverage, veyltrace:off} with a mean set size of \resnum{re2-oracle-rq3-v1}{mean set size, veyltrace:off}.

The rest of the paper is organised as follows.
Section~\ref{sec:related} places \sys among statistical, causal and LLM-based RCA methods.
Section~\ref{sec:design} describes the design, Section~\ref{sec:eval} the pre-registered evaluation and Section~\ref{sec:results} the results.
Section~\ref{sec:discussion} discusses what the verifier and the LLM each contribute, and Section~\ref{sec:threats} the threats to validity.
